\section{Before the modern taxonomy: Scripture, Fathers, and councils}

The New Testament already shows authoritative transmission: the apostolic
Gospel is received and handed on (1 Cor 15:1--5), the ``pattern of sound words''
is guarded (2 Tim 1:13--14), and the Jerusalem assembly distinguishes the
Gospel's demand from an unnecessary burden (Acts 15). These texts do not supply
a ready-made 1998 chart. They disclose the apostolic realities from which the
Church's later precision grows.

Irenaeus answers Gnostic appeals to secret knowledge with the public apostolic
tradition confessed throughout the churches (\emph{Against Heresies}
1.10; 3.3--4). He joins unity of the received faith to legitimate growth in
explanation: greater skill can draw out connections without changing the
faith's subject matter. His appeal to public succession likewise rejects a
private elite's claim to possess a second apostolic teaching. This is real
patristic provenance for creed, reception, and teaching office; it is not yet
the modern three-paragraph taxonomy.

Ancient creeds and councils arose from baptismal confession, worship, exegesis,
and controversy. Nicaea's \emph{homoousios} did not add a second revelation; it
made a judgment about which account faithfully confessed the scriptural Son.
Ephesus and Chalcedon likewise used precise exclusions to protect the identity
of Christ. Reception was real and sometimes contested, but ``reception''
cannot mean that a definition becomes true only when each believer approves
it. Nor should later categories be projected backwards as though every Father
had assigned a modern theological note.

Medieval schools developed increasingly exact censures and ``theological
notes'' to distinguish revealed doctrine, proximate conclusions, common
teaching, probable opinion, and error. These systems did valuable diagnostic
work but never formed one universally promulgated, immutable table. Labels
such as \emph{de fide}, \emph{fidei proxima}, \emph{sententia certa}, and
\emph{communis} must be attributed to an author and date. They are historical
theological judgments, not substitutes for the postconciliar profession and
the competent Magisterium.

